\documentclass[a4paper, 10pt]{article}
\usepackage[utf8]{inputenc} % Change according your file encoding
\usepackage{graphicx}
\usepackage{url}

%opening
\title{Seminar Report: Paxy}
\author{\textbf{Name of team members}}
\date{\normalsize\today{}}

\begin{document}

\maketitle

\begin{center}
Upload your report in PDF format.
  
Use this LaTeX template to format the report.
  
Validate your basic Paxy code through the review questionnaires at ATENEA.

Submit a compressed file (.tar.gz) containing ALL your source code files.
\end{center}

\section{Introduction}

\textit{Introduce in a couple of sentences the seminar and the main topic related to distributed systems it covers.}

\section{Experiments and Questions}

\subsection{Message delays}

\begin{table}[htbp]
\centering
\small
\begin{tabular}{r|r|r|r|r}
\hline
Delay (ms) & Rounds & Promise & Vote & Both \\
\hline
0 & 1 & 8 & 8 & 8 \\
500 & 1 & 704 & 509 & 1204 \\
1000 & 1 & 1204 & 1009 & 2205 \\
1500 & 1 & 1704 & 1509 & 3205 \\
2000 & --- & --- & --- & --- \\
\hline
\end{tabular}
\caption{Consensus time (ms) and rounds needed to reach consensus (1 = first ballot). \texttt{---} means no decision within 12~s. Timeout 2000~ms.}
\label{tab:delays}
\end{table}

\noindent Delaying only votes adds about the delayed time to the consensus time, because promises are immediate, so the first proposer wins. Delaying only promises adds about the delayed time $+200$~ms. We can see in the table that the relation between delay and consensus time seems linear. Also we can see that delay doesn't affect to the number of rounds, but this happens just in case the delay remains less than the timeout, it just affects to the consensus time. However if the delay is not constant and sometimes votes or promises experience delay times larger than the timeout then it will cost more rounds. \\

\noindent Increasing the timeout does not reduce the consensus time if the delays are less than the timeout, but it slows down failure detection and makes losing proposers wait longer. if the system is asynchronous no timeout guarantees termination.

\subsection{\texttt{sorry} messages}

\textit{\textbf{Experiment.} Avoid sending \texttt{sorry} messages by commenting the corresponding sentences in the acceptor. Compare the time and the number of rounds needed to get consensus with and without \texttt{sorry} messages.} 

{\setlength{\parskip}{10pt}\noindent\textit{\textbf{Question.} Discuss the impact of dropping \texttt{sorry} messages, particularly, whether you could still come to an agreement.}}

\subsection{Lost messages}

\textit{\textbf{Experiment.} Try randomly dropping \texttt{promise} messages, \texttt{vote} messages, or both in the acceptor. Provide figures or tables with experimental results comparing the time and the number of rounds needed to get consensus.}
 
{\setlength{\parskip}{10pt}\noindent\textit{\textbf{Question.} Discuss the impact of dropping messages, and indicate the percentage of messages we can drop until consensus is no longer possible.}}

\subsection{Number of \textit{acceptors} and \textit{proposers}}

We increased the number of acceptors (3, 5, 7, 9) with 3 proposers, and the
number of proposers (1, 3, 5, 7) with 5 acceptors. All proposers start
together, with no message delays and the default timeout (2000~ms); each point
is the mean of 10 headless runs (Tables~\ref{tab:acceptors}
and~\ref{tab:proposers}).

\begin{table}[htbp]
\centering
\small
\begin{tabular}{c|c|c|c|c}
\hline
Acceptors & $f$ & $T_{\mathrm{cons}}$ (ms) & Winner/max rounds & Messages \\
\hline
3 & 1 & 16 & 1 / 2.6 & 52 \\
5 & 2 & 16 & 1 / 2.6 & 86 \\
7 & 3 & 16 & 1 / 2.7 & 122 \\
9 & 4 & 16 & 1 / 2.7 & 157 \\
\hline
\end{tabular}
\caption{Number of acceptors (mean of 10 runs, 3 proposers, timeout 2000~ms).
$f=(A-1)/2$ is the number of crashed acceptors tolerated; rounds are
winner/maximum; messages are the total per run.}
\label{tab:acceptors}
\end{table}

\begin{table}[htbp]
\centering
\small
\begin{tabular}{c|c|c|c|c}
\hline
Proposers & $T_{\mathrm{cons}}$ (ms) & $T_{\mathrm{all}}$ (ms) & Winner/max rounds & Messages \\
\hline
1 & 16 & 16 & 1 / 1 & 20 \\
3 & 16 & 3635 & 1 / 2.8 & 88 \\
5 & 16 & 5057 & 1 / 3.5 & 163 \\
7 & 16 & 6687 & 1 / 4.3 & 256 \\
\hline
\end{tabular}
\caption{Number of proposers (mean of 10 runs, 5 acceptors, timeout 2000~ms).
$T_{\mathrm{all}}$ is the time until all proposers decide; rounds are
winner/maximum; messages are the total per run.}
\label{tab:proposers}
\end{table}

With more acceptors the quorum grows ($\lfloor A/2\rfloor+1$), so every ballot
exchanges more messages and the system tolerates $f=(A-1)/2$ crashed
acceptors, but the consensus time and the winning round stay roughly constant.
The price is message traffic and a larger quorum to assemble.

With more proposers, consensus is still reached by one proposer in its first
ballot. The others lose the race, receive \texttt{sorry}, and retry, so the
maximum number of rounds, the time until all decide, and the traffic grow with
the number of proposers. More proposers therefore improve availability (any
one can drive consensus) but add contention and load; with message delays
(Section~2.1) this contention becomes worse.

Each ballot exchanges about $4A$ messages ($A$ prepares and $A$ accepts, plus
their replies), so traffic grows linearly with the number of acceptors and
with the number of retries.

\subsection{Distributed execution}

\textit{\textbf{Code.} Discuss the changes needed in the \texttt{paxy} module to create the proposers in a remote Erlang instance (named \texttt{paxy-pro}) and to ensure that they can connect correctly to the acceptors, which must be created in a different remote Erlang instance (named \texttt{paxy-acc}).}

{\setlength{\parskip}{10pt}\noindent\textit{\textbf{Experiment.} Record a video providing empirical evidence of the distributed execution working properly.}}

\subsection{Fault tolerance}

\textit{\textbf{Code.} Discuss the changes needed in the acceptor to make it fault tolerant.}

{\setlength{\parskip}{10pt}\noindent\textit{\textbf{Experiment.} Record a video demonstrating that a crashed acceptor can recover successfully with exactly the same state it had before failing.}}

\subsection{Improvement based on \texttt{sorry} messages}

\textit{\textbf{Code.} Discuss the changes needed in the \texttt{collect/4} and \texttt{vote/2} procedures in the proposer.}
 
{\setlength{\parskip}{10pt}\noindent\textit{\textbf{Experiment.} Compare the time needed to get consensus with and without the improvement.}}

\section{Personal opinion}

\textit{Provide your personal opinion of the seminar, indicating whether it should be included in next year's course or not.}

\end{document}
